\section{Discrepancies across one flip}\label{sec:comparison}

We first establish the properties of an individual step that will be
used throughout the termination argument.  The two relative ample
signs produce an effective divisor on the normalized graph.  Its
support contains the entire inverse image of the exceptional loci;
this is what gives strict increase for every place centred in either
locus.

\begin{proposition}[Comparison across a flip]\label{prop:comparison}
Let
\[
 T\overset{f}{\longrightarrow}Z
   \overset{g}{\longleftarrow}T^+
\]
be a step satisfying the hypotheses of Theorem~\ref{thm:termination} between normal compact K\"ahler fourfolds.  Write
\[
 D_T=K_T+B_T+M_T,\qquad
 D_{T^+}=K_{T^+}+B_{T^+}+M_{T^+},
\]
where the nef traces come from the same fixed b-divisor $\bM$.
Then the following assertions hold.
\begin{enumerate}
\item The exceptional images coincide:
\[
 A:=f(\Exc f)=g(\Exc g).
\]
With reduced structures, put
\[
 L=f^{-1}(A)=\Exc f,\qquad
 L^+=g^{-1}(A)=\Exc g.
\]
These are closed analytic subsets, and
\[
 T\setminus L\simeq T^+\setminus L^+,
 \qquad \dim L,\dim L^+\leq2,\qquad \dim A\leq1.
\]
\item For every divisorial place $E$ over the two models,
\begin{equation}\label{eq:discrepancy-monotonicity}
 a(E;T^+,B_{T^+}+\bM)
 \geq a(E;T,B_T+\bM).
\end{equation}
The inequality is strict if the centre of $E$ on $T$ is contained in
$L$, or if its centre on $T^+$ is contained in $L^+$.
A place is exceptional over $T$ if and only if it is exceptional over
$T^+$. In particular, the generalized canonical condition of
Theorem~\ref{thm:termination} is preserved. If it holds on $T$, then
every exceptional place with centre contained in $L^+$ satisfies
\begin{equation}\label{eq:target-strict-discrepancy}
 a(E;T^+,B_{T^+}+\bM)>1.
\end{equation}
\item The exceptional loci satisfy
\begin{equation}\label{eq:exceptional-dimensions}
 \dim L+\dim L^+\geq3.
\end{equation}
Consequently, if $L^+$ contains no surface, then $L$ contains a surface.
\end{enumerate}
\end{proposition}

\begin{proof}
\emph{The common exceptional image.}
For each of the two morphisms, the exceptional locus is precisely the
inverse image of the locus of positive-dimensional fibres, by the
discussion in Section~\ref{subsec:places}. Off the union of these
two images, the morphisms initially give isomorphic opens whose
complements upstairs have codimension at least two. Thus prime
divisors already correspond, before equality of the images is known.

Let $U\subset Z$ be the open set where $f$ is an isomorphism.  Transport
$D_T|_{f^{-1}(U)}$ to an actual $\Q$-Cartier divisor $D_U$ on $U$.
The two adjoints agree under the identification of prime divisors
induced by the small map: this holds for the canonical and boundary
terms by their prescribed transforms, and for the nef traces by the
pushforward rule for the fixed b-divisor.  Hence on $g^{-1}(U)$ the
divisors $D_{T^+}$ and $g^*D_U$ agree away from the exceptional locus
of $g$.  That locus contains no divisor, so the two Weil divisors
agree everywhere on $g^{-1}(U)$.  Clearing Cartier multiples and using
normality gives equality as $\Q$-Cartier divisors as well.

If $g$ had a positive-dimensional fibre over a point of $U$, its
projectivity would provide a curve $C$ in that fibre.  The equality
just proved would give $D_{T^+}\cdot C=0$, contrary to $g$-ampleness.
Thus $g$ is also an isomorphism over $U$.  Interchanging $f$ and $g$
and using $f$-ampleness of $-D_T$ proves the reverse inclusion of their
isomorphism loci.  This establishes the asserted common image $A$
and the identification of the complements of $L$ and $L^+$.
Smallness gives $\dim L,\dim L^+\leq2$.  Every fibre over $A$ is
positive dimensional, so the fibre-dimension theorem gives
$\dim A\leq1$.

\emph{An effective divisor on the graph.}
Let $G$ be the normalization of the component of $T\times_ZT^+$
dominating $Z$.  Its projections fit into the commutative diagram
\begin{equation}\label{eq:normalized-graph}
 \begin{array}{ccc}
 G&\xrightarrow{\ q\ }&T^+\\[3pt]
 {\scriptstyle p}\big\downarrow&&
                  \big\downarrow{\scriptstyle g}\\[3pt]
 T&\xrightarrow{\ f\ }&Z.
 \end{array}
 \qquad h:=f\circ p=g\circ q.
\end{equation}
The morphisms $p,q,h$ are projective and bimeromorphic, as in
Proposition~\ref{prop:common-models}.  Define the actual $\Q$-Cartier
divisor
\begin{equation}\label{eq:graph-divisor}
 F:=p^*D_T-q^*D_{T^+}.
\end{equation}
It vanishes off $h^{-1}(A)$.  Since both $L$ and $L^+$ have codimension
at least two, $F$ is exceptional over both models; in particular,
$p_*F=0$.

If $C$ is a curve contracted by $p$, then $q_*C$ is either zero or a
curve contracted by $g$.  The projection formula therefore gives
\[
 (-F)\cdot C=D_{T^+}\cdot q_*C\geq0.
\]
Thus $-F$ is $p$-nef in the relative curve-degree sense.  Applying
Lemma~\ref{lem:negativity} to a positive Cartier multiple of $F$
proves that $F\geq0$.

There is a strict inequality on every curve contracted by $h$.
Indeed, for such an irreducible curve $C$, the relative ample signs
give
\[
 D_T\cdot p_*C\leq0,\qquad
 D_{T^+}\cdot q_*C\geq0,
\]
with strict inequality whenever the corresponding image is a curve.
The two images cannot both be points: the map
$G\to T\times_ZT^+$ is finite onto the graph component, so it cannot
contract $C$.  Hence
\begin{equation}\label{eq:graph-curve-negativity}
 F\cdot C=D_T\cdot p_*C-D_{T^+}\cdot q_*C<0
 \qquad\text{for every $h$-contracted irreducible curve $C$.}
\end{equation}

We next show that this strict inequality determines the whole support:
\begin{equation}\label{eq:graph-support}
 \Supp F=h^{-1}(A).
\end{equation}
Fix $z\in A$.  The fibre $h^{-1}(z)$ is projective and connected;
connectedness also follows directly from the proper bimeromorphic
morphism $h$ and normality of $Z$.  It is positive dimensional because
it surjects onto $f^{-1}(z)$.  Its reduction has finitely many
irreducible components.  A zero-dimensional irreducible component
would be a point disjoint from every other component, contradicting
connectedness.  Therefore every point belongs to a positive-dimensional
projective component.  Intersecting such a component with hyperplanes
through the point produces an irreducible curve through that point.

An effective $\Q$-Cartier divisor has nonnegative degree on every
compact irreducible curve not contained in its support: a positive
multiple is an effective Cartier divisor, whose canonical section
restricts to a nonzero holomorphic section on the normalization of
that curve.  By \eqref{eq:graph-curve-negativity}, every curve in
$h^{-1}(z)$ is therefore contained in $\Supp F$.  The preceding
paragraph puts every point of the fibre in $\Supp F$.  This proves
$h^{-1}(A)\subset\Supp F$; the reverse inclusion was already noted.

\emph{Discrepancies of all places.}
Choose a sufficiently high smooth common model carrying the fixed
nef data and dominating $G$, as supplied by
Proposition~\ref{prop:common-models}.  To compare a specified place
$E$, further dominate a model representing it, so that $E$ is a prime
divisor on a smooth model $W$ with morphism $r:W\to G$.  The same
canonical divisor $K_W$ and the same nef trace $M_W$ occur in both
discrepancy equations:
\[
 \begin{aligned}
 K_W+\Delta_T+M_W&=(p\circ r)^*D_T,\\
 K_W+\Delta_{T^+}+M_W&=(q\circ r)^*D_{T^+}.
 \end{aligned}
\]
Their difference is the equality of actual divisors
\[
 r^*F=\Delta_T-\Delta_{T^+}.
\]
Taking the coefficient at $E$ yields
\begin{equation}\label{eq:graph-discrepancy-identity}
 a(E;T^+,B_{T^+}+\bM)-a(E;T,B_T+\bM)
   =\operatorname{coeff}_E(r^*F).
\end{equation}
Pullback of the effective $\Q$-Cartier divisor $F$ is effective, proving
\eqref{eq:discrepancy-monotonicity}.

If the centre of $E$ on $T$ lies in $L$, its centre on $G$ lies in
$p^{-1}(L)=h^{-1}(A)$; the same conclusion holds if its centre on
$T^+$ lies in $L^+$.  By \eqref{eq:graph-support}, this centre is
contained in $\Supp F$.  A local equation of a positive effective
Cartier multiple of $F$ vanishes on that centre.  Its pullback has
positive order along $E$: it is a nonzero holomorphic function on
the inverse image of the coordinate neighbourhood and vanishes on
the part of $E$ lying there.  Consequently the
right-hand side of \eqref{eq:graph-discrepancy-identity} is strictly
positive.  This proves strictness for an arbitrary place with either
specified centre condition.

A place whose centre is a prime divisor on a normal model is
represented by that divisor, since a proper bimeromorphic morphism
to a normal space is an isomorphism at general points of each prime
divisor on the target.  Every prime divisor
on $T$ meets $T\setminus L$, and its corresponding prime divisor on
$T^+$ meets $T^+\setminus L^+$.  The common isomorphism identifies
their orders. Thus a place is nonexceptional on one side exactly
when it is nonexceptional on the other. Together with
\eqref{eq:discrepancy-monotonicity}, this preserves positivity for
all places and the lower bound $1$ for exceptional places. For a
place centred in $L^+$, strictness improves that lower bound to
\eqref{eq:target-strict-discrepancy}.

\emph{The dimension inequality.}
The step is nontrivial, so $A\neq\varnothing$.  Equation
\eqref{eq:graph-support} implies $F\neq0$.  Every irreducible component
$P$ of $\Supp F$ is a divisor on the four-dimensional normal space
$G$, hence has dimension three.  Its finite image under
$G\to T\times_ZT^+$ is contained in $L\times_ZL^+$.  Therefore
\[
 3=\dim P\leq\dim(L\times_ZL^+)\leq\dim L+\dim L^+,
\]
which proves \eqref{eq:exceptional-dimensions}.  If $L^+$ contains no
surface, then $\dim L^+\leq1$, while smallness gives $\dim L\leq2$.
The inequality forces $\dim L=2$, so a two-dimensional irreducible
component of $L$ is the required surface.
\end{proof}

We will call an irreducible compact analytic surface contained in $L$
a \emph{flipping surface}, and one contained in $L^+$ a
\emph{flipped surface}.  Applying Proposition~\ref{prop:comparison}
successively shows that every pair in a sequence as in
Theorem~\ref{thm:termination} remains generalized canonical. The
strict bound \eqref{eq:target-strict-discrepancy} concerns places
centred in the flipped locus and will supply the local terminal
condition used to study its surfaces.
